\documentclass[a4paper,11pt]{article}
\usepackage[utf8]{inputenc}
\usepackage[T1]{fontenc}
\usepackage[ngerman]{babel}
\usepackage{geometry}
\geometry{left=2.5cm, right=2.5cm, top=2.5cm, bottom=2.5cm}
\usepackage{longtable}
\usepackage{booktabs}
\usepackage{array}
\usepackage{xcolor}
\usepackage{colortbl}
\usepackage{hyperref}
\usepackage{fancyhdr}
\usepackage{tabularx}
\usepackage{enumitem}
\usepackage{graphicx}
\usepackage{amssymb}

\definecolor{headerblue}{HTML}{1565C0}
\definecolor{tableheader}{HTML}{E3F2FD}
\definecolor{pass}{HTML}{2E7D32}
\definecolor{warn}{HTML}{F57F17}
\definecolor{fail}{HTML}{D32F2F}

\hypersetup{
    colorlinks=true,
    linkcolor=headerblue,
    urlcolor=headerblue
}

\pagestyle{fancy}
\fancyhf{}
\fancyhead[L]{\textbf{INSPECTAIR}}
\fancyhead[R]{Testspezifikation v2.0}
\fancyfoot[C]{\thepage}
\renewcommand{\headrulewidth}{0.4pt}

\setlength{\parindent}{0pt}
\setlength{\parskip}{6pt}

\begin{document}

\begin{center}
{\LARGE\textbf{INSPECTAIR}}\\[0.3cm]
{\Large Testspezifikation v2.0}\\[0.3cm]
{\large DHBW Ravensburg \textbar{} Mikrocomputertechnik 1}\\[0.5cm]
\begin{tabular}{ll}
\textbf{Projekt:} & InspectAir -- Luftqualitätsmessgerät \\
\textbf{Version:} & 2.0 \\
\textbf{Datum:} & Februar 2026 \\
\textbf{Autoren:} & Maximilian Liebl, Cedric Stroh, Jannis Messer \\
\textbf{Modul:} & Mikrocomputertechnik 1, DHBW Ravensburg \\
\textbf{Dozent:} & Thomas Stingl (Airbus) \\
\end{tabular}
\end{center}

\vspace{0.5cm}
\noindent\textit{Dieses Dokument wurde unter Verwendung von KI-Tools (Claude, Anthropic) zur Überprüfung erstellt.}
\vspace{0.3cm}

\tableofcontents
\newpage

% ============================================================
\section{Übersicht}
% ============================================================

Diese Testspezifikation enthält \textbf{17 Testfälle}, die alle Features des InspectAir-Systems abdecken.

\begin{longtable}{p{5.5cm} p{2cm} p{2.5cm} p{2.5cm}}
\toprule
\rowcolor{tableheader}
\textbf{Kategorie} & \textbf{Anzahl} & \textbf{Bestanden} & \textbf{Aufwand} \\
\midrule
\endhead
Einzeltests Sensoren                        & 6 & 6 \textcolor{pass}{\checkmark} & erledigt \\
Hardware-Tests (MOSFET, Kondensatoren)      & 2 & 1 \textcolor{warn}{$\circlearrowleft$} & \textasciitilde30\,Min \\
Integrationstest                            & 1 & 1 \textcolor{warn}{$\circlearrowleft$} & \textasciitilde15\,Min \\
UI-Tests (5 Screens, Button, Konsistenz)    & 4 & 0 & \textasciitilde45\,Min \\
Power Management (Dimming, Sleep)           & 3 & 0 & \textasciitilde30\,Min \\
Systemtest                                  & 1 & 0 & \textasciitilde30\,Min \\
\midrule
\textbf{GESAMT} & \textbf{17} & \textbf{6/17} & \textbf{\textasciitilde2.5 Stunden} \\
\bottomrule
\end{longtable}

\subsection{Getestete Features}

\begin{longtable}{p{6.5cm} p{6.5cm}}
\toprule
\rowcolor{tableheader}
\textbf{Feature} & \textbf{Relevante Tests} \\
\midrule
\endhead
BS170 MOSFET Backlight-Dimming                 & TC-07, TC-14, TC-15 \\
Stützkondensatoren (220\,µF 25V)               & TC-08, TC-09 \\
5 UI-Screens (Tree, Overview, Detail, Analog, Bubble) & TC-10, TC-11, TC-12, TC-13 \\
Radar $\rightarrow$ Display aufwecken          & TC-15, TC-16, TC-17 \\
Display dimmen nach Timeout                     & TC-14, TC-17 \\
ESP32 Light Sleep Modus                        & TC-16 \\
\bottomrule
\end{longtable}

% ============================================================
\section{Teil 1: Einzeltests Sensoren (bestanden)}
% ============================================================

Diese Tests wurden bereits bei der Komponentenbeschaffung durchgeführt.

\begin{longtable}{p{3cm} p{10cm}}
\toprule
\rowcolor{tableheader}
\multicolumn{2}{l}{\textbf{TC-01 -- Einzeltest: Display ST7796S \textcolor{pass}{\checkmark}}} \\
\midrule
\endhead
\textbf{Beschreibung}   & SPI-Kommunikation und Farbdarstellung des 4-Zoll-Displays. \\
\textbf{Vorbedingung}    & Display verkabelt: CS$\rightarrow$GPIO\,14, DC$\rightarrow$GPIO\,9, MOSI$\rightarrow$GPIO\,11, CLK$\rightarrow$GPIO\,12, RST$\rightarrow$GPIO\,46. \\
\textbf{Testschritte}    & 1. Testfirmware flashen \newline 2. RGB-Vollbildfarben prüfen (Rot, Grün, Blau) \newline 3. Text-Rendering testen \\
\textbf{Pass-Kriterium}  & Display zeigt alle Farben ohne Artefakte. Text scharf lesbar. \\
\textbf{Ergebnis}        & \textcolor{pass}{$\boxtimes$ PASS} $\square$ FAIL \quad Datum: Jan 2026 \quad Tester: Team \\
\bottomrule
\end{longtable}

\begin{longtable}{p{3cm} p{10cm}}
\toprule
\rowcolor{tableheader}
\multicolumn{2}{l}{\textbf{TC-02 -- Einzeltest: AHT20 (Temp/Feuchte) \textcolor{pass}{\checkmark}}} \\
\midrule
\endhead
\textbf{Beschreibung}   & I2C-Kommunikation und plausible Messwerte. \\
\textbf{Vorbedingung}    & AHT20 an I2C: SDA$\rightarrow$GPIO\,8, SCL$\rightarrow$GPIO\,18. Pull-ups 4.7\,k$\Omega$. \\
\textbf{Testschritte}    & 1. I2C-Scan (Adresse 0x38) \newline 2. Temperatur und Feuchte auslesen \\
\textbf{Pass-Kriterium}  & Werte: 15--30\,°C, 30--70\,\% rH. \\
\textbf{Ergebnis}        & \textcolor{pass}{$\boxtimes$ PASS} $\square$ FAIL \quad Datum: Jan 2026 \quad Tester: Team \\
\bottomrule
\end{longtable}

\begin{longtable}{p{3cm} p{10cm}}
\toprule
\rowcolor{tableheader}
\multicolumn{2}{l}{\textbf{TC-03 -- Einzeltest: SGP40 (VOC) \textcolor{pass}{\checkmark}}} \\
\midrule
\endhead
\textbf{Beschreibung}   & I2C-Kommunikation und VOC-Rohwerte. \\
\textbf{Vorbedingung}    & SGP40 an I2C parallel zu AHT20. Adresse 0x59. \\
\textbf{Testschritte}    & 1. I2C-Scan (Adresse 0x59) \newline 2. SRAW-Wert nach 60\,s Aufwärmzeit lesen \\
\textbf{Pass-Kriterium}  & SRAW: 15000--35000. Selbsttest OK. \\
\textbf{Ergebnis}        & \textcolor{pass}{$\boxtimes$ PASS} $\square$ FAIL \quad Datum: Jan 2026 \quad Tester: Team \\
\bottomrule
\end{longtable}

\begin{longtable}{p{3cm} p{10cm}}
\toprule
\rowcolor{tableheader}
\multicolumn{2}{l}{\textbf{TC-04 -- Einzeltest: MH-Z19C (CO\textsubscript{2}) \textcolor{pass}{\checkmark}}} \\
\midrule
\endhead
\textbf{Beschreibung}   & UART-Kommunikation und CO\textsubscript{2}-Messung. \\
\textbf{Vorbedingung}    & MH-Z19C: TX$\rightarrow$GPIO\,4, RX$\leftarrow$GPIO\,5 (SoftwareSerial). VCC$\rightarrow$5V. Stützkondensator 220\,µF (25V) am VCC! \\
\textbf{Testschritte}    & 1. SoftwareSerial 9600 Baud initialisieren \newline 2. 3\,Min Aufwärmzeit \newline 3. CO\textsubscript{2}-Wert bei Frischluft messen \\
\textbf{Pass-Kriterium}  & CO\textsubscript{2}: 350--600\,ppm bei Frischluft. Keine Timeouts. \\
\textbf{Ergebnis}        & \textcolor{pass}{$\boxtimes$ PASS} $\square$ FAIL \quad Datum: Jan 2026 \quad Tester: Team \\
\bottomrule
\end{longtable}

\begin{longtable}{p{3cm} p{10cm}}
\toprule
\rowcolor{tableheader}
\multicolumn{2}{l}{\textbf{TC-05 -- Einzeltest: PMS5003 (Feinstaub) \textcolor{pass}{\checkmark}}} \\
\midrule
\endhead
\textbf{Beschreibung}   & UART-Kommunikation und PM-Werte. \\
\textbf{Vorbedingung}    & PMS5003: TX$\rightarrow$GPIO\,16, RX$\leftarrow$GPIO\,17. VCC$\rightarrow$5V. UART1. \\
\textbf{Testschritte}    & 1. UART 9600 Baud \newline 2. Auf Datenframes warten \newline 3. PM2.5-Wert prüfen \\
\textbf{Pass-Kriterium}  & Gültige Frames mit korrekter Checksum. PM2.5: 0--100\,µg/m³. \\
\textbf{Ergebnis}        & \textcolor{pass}{$\boxtimes$ PASS} $\square$ FAIL \quad Datum: Jan 2026 \quad Tester: Team \\
\bottomrule
\end{longtable}

\begin{longtable}{p{3cm} p{10cm}}
\toprule
\rowcolor{tableheader}
\multicolumn{2}{l}{\textbf{TC-06 -- Einzeltest: LD2410C (Radar) \textcolor{pass}{\checkmark}}} \\
\midrule
\endhead
\textbf{Beschreibung}   & UART + Level Shifter, Präsenzerkennung. \\
\textbf{Vorbedingung}    & LD2410C über Level Shifter: TX$\rightarrow$GPIO\,7, RX$\leftarrow$GPIO\,6, OUT$\rightarrow$GPIO\,15. Stützkondensator 220\,µF (25V)! \\
\textbf{Testschritte}    & 1. Level Shifter TX-Ausgang messen (MUSS 3.3V sein!) \newline 2. UART 256000 Baud \newline 3. Präsenz bei Annäherung testen \\
\textbf{Pass-Kriterium}  & TX-Pegel $\leq$3.3V. OUT-Pin schaltet bei Bewegung. \\
\textbf{Ergebnis}        & \textcolor{pass}{$\boxtimes$ PASS} $\square$ FAIL \quad Datum: Jan 2026 \quad Tester: Team \\
\bottomrule
\end{longtable}

% ============================================================
\section{Teil 2: Hardware-Tests}
% ============================================================

MOSFET-Steuerung und Kondensator-Stabilisierung.

\begin{longtable}{p{3cm} p{10cm}}
\toprule
\rowcolor{tableheader}
\multicolumn{2}{l}{\textbf{TC-07 -- Hardware: BS170 MOSFET Backlight-Steuerung}} \\
\midrule
\endhead
\textbf{Beschreibung}   & Testet die PWM-Dimmung des Display-Backlights über den BS170 MOSFET. \\
\textbf{Vorbedingung}    & BS170 MOSFET: Gate$\rightarrow$GPIO\,3 (über 1\,k$\Omega$), Drain$\rightarrow$Display BL-, Source$\rightarrow$GND. \\
\textbf{Testschritte}    & 1. PWM-Kanal auf GPIO\,3 konfigurieren (44.1\,kHz) \newline 2. Duty 0\,\% setzen $\rightarrow$ Display komplett dunkel \newline 3. Duty 50\,\% setzen $\rightarrow$ halbe Helligkeit \newline 4. Duty 100\,\% setzen $\rightarrow$ volle Helligkeit \newline 5. Rampe 0$\rightarrow$100\,\% in 1\,s (Fade-In Effekt) \\
\textbf{Pass-Kriterium}  & Helligkeit stufenlos regelbar. Kein Flackern. Smooth Fade möglich. \\
\textbf{Ergebnis}        & $\square$ PASS \quad $\square$ FAIL \quad Datum: \_\_\_\_\_\_ \quad Tester: \_\_\_\_\_\_ \\
\bottomrule
\end{longtable}

\begin{longtable}{p{3cm} p{10cm}}
\toprule
\rowcolor{tableheader}
\multicolumn{2}{l}{\textbf{TC-08 -- Hardware: Stützkondensatoren Stabilität \textcolor{warn}{$\circlearrowleft$}}} \\
\midrule
\endhead
\textbf{Beschreibung}   & Prüft, ob die 220\,µF (25V) Kondensatoren am CO\textsubscript{2}-Sensor und Radar Stromspitzen puffern und Abstürze verhindern. \\
\textbf{Vorbedingung}    & 220\,µF (25V) Elkos an VCC/GND von MH-Z19C und LD2410C. Alle Sensoren aktiv. \\
\textbf{Testschritte}    & 1. System OHNE Kondensatoren starten $\rightarrow$ Absturz provozieren \newline 2. Kondensatoren einlöten \newline 3. System 10 Minuten laufen lassen \newline 4. Spannung am 5V-Rail mit Oszilloskop/Multimeter beobachten \\
\textbf{Pass-Kriterium}  & MIT Kondensatoren: Keine Abstürze, 5V-Rail stabil (4.8--5.2V). OHNE: Dokumentierter Absturz. \\
\textbf{Ergebnis}        & $\square$ PASS \quad $\square$ FAIL \quad Datum: \_\_\_\_\_\_ \quad Tester: \_\_\_\_\_\_ \\
\bottomrule
\end{longtable}

% ============================================================
\section{Teil 3: Integrationstest}
% ============================================================

\begin{longtable}{p{3cm} p{10cm}}
\toprule
\rowcolor{tableheader}
\multicolumn{2}{l}{\textbf{TC-09 -- Integration: Alle Sensoren gleichzeitig \textcolor{warn}{$\circlearrowleft$}}} \\
\midrule
\endhead
\textbf{Beschreibung}   & Prüft, ob alle 5 Sensoren + Radar parallel funktionieren ohne Konflikte. \\
\textbf{Vorbedingung}    & Alle Sensoren verkabelt. Stützkondensatoren installiert. Firmware mit allen Treibern. \\
\textbf{Testschritte}    & 1. System starten \newline 2. Serial-Log auf Init-Meldungen prüfen (alle [OK]) \newline 3. 5 Minuten alle Werte loggen \newline 4. Auf I2C/UART-Fehler oder Timeouts prüfen \\
\textbf{Pass-Kriterium}  & Alle Sensoren initialisiert. Fehlerrate $<$1\,\% über 5 Minuten. \\
\textbf{Ergebnis}        & $\square$ PASS \quad $\square$ FAIL \quad Datum: \_\_\_\_\_\_ \quad Tester: \_\_\_\_\_\_ \\
\bottomrule
\end{longtable}

% ============================================================
\section{Teil 4: UI-Tests}
% ============================================================

Tests für alle 5 UI-Screens und die Umschaltung per Button.

Die 5 Screens sind:
\begin{enumerate}[leftmargin=2em]
    \item \textbf{Tree-Animation} -- Startbildschirm mit animiertem Baum
    \item \textbf{Overview} -- Große AQI-Anzeige + 2 Kacheln (minimalistisch)
    \item \textbf{Detail} -- Kleine AQI + 4 Kacheln (volle Informationen)
    \item \textbf{Analog} -- Analoge Cockpit-Instrumente
    \item \textbf{Bubble} -- Dynamische Kreise (Modern Bubbles)
\end{enumerate}

\begin{longtable}{p{3cm} p{10cm}}
\toprule
\rowcolor{tableheader}
\multicolumn{2}{l}{\textbf{TC-10 -- UI: Screen-Darstellung (alle 5 Screens)}} \\
\midrule
\endhead
\textbf{Beschreibung}   & Testet die korrekte Darstellung aller 5 UI-Screens mit allen Messwerten. \\
\textbf{Vorbedingung}    & Display und alle Sensoren funktional. Alle 5 Screens implementiert (LVGL 9.2). \\
\textbf{Testschritte}    & 1. System starten (Standard: Tree-Animation) \newline 2. Zu Overview-Screen wechseln -- alle Werte prüfen \newline 3. Zu Detail-Screen wechseln -- alle 6 Werte prüfen: CO\textsubscript{2}, PM2.5, VOC, Temp, Feuchte, Präsenz \newline 4. Zu Analog-Screen wechseln -- Instrumentendarstellung prüfen \newline 5. Zu Bubble-Screen wechseln -- Kreisdarstellung prüfen \newline 6. Farbcodierung prüfen (Grün/Gelb/Rot je nach Grenzwert) \newline 7. Lesbarkeit aus 1\,m Entfernung \\
\textbf{Pass-Kriterium}  & Alle Werte auf allen Screens sichtbar und korrekt. Farben entsprechen Grenzwerten. Gut lesbar. \\
\textbf{Ergebnis}        & $\square$ PASS \quad $\square$ FAIL \quad Datum: \_\_\_\_\_\_ \quad Tester: \_\_\_\_\_\_ \\
\bottomrule
\end{longtable}

\begin{longtable}{p{3cm} p{10cm}}
\toprule
\rowcolor{tableheader}
\multicolumn{2}{l}{\textbf{TC-11 -- UI: Screen-spezifische Features}} \\
\midrule
\endhead
\textbf{Beschreibung}   & Testet die individuellen Features jedes Screens. \\
\textbf{Vorbedingung}    & Alle 5 Screens implementiert. \\
\textbf{Testschritte}    & 1. Tree-Screen: Animation läuft flüssig, reagiert auf Luftqualität \newline 2. Overview-Screen: Große AQI-Zahl korrekt, 2 Kacheln lesbar \newline 3. Detail-Screen: Alle Werte + Zusatzinfos korrekt, 4 Kacheln \newline 4. Analog-Screen: Zeiger bewegen sich entsprechend der Messwerte \newline 5. Bubble-Screen: Kreise ändern Größe/Farbe dynamisch \\
\textbf{Pass-Kriterium}  & Alle screen-spezifischen Features funktional. Design konsistent. Keine Rendering-Artefakte. \\
\textbf{Ergebnis}        & $\square$ PASS \quad $\square$ FAIL \quad Datum: \_\_\_\_\_\_ \quad Tester: \_\_\_\_\_\_ \\
\bottomrule
\end{longtable}

\begin{longtable}{p{3cm} p{10cm}}
\toprule
\rowcolor{tableheader}
\multicolumn{2}{l}{\textbf{TC-12 -- UI: Button-Wechsel zwischen Screens}} \\
\midrule
\endhead
\textbf{Beschreibung}   & Testet den zyklischen Wechsel zwischen allen 5 Screens per Button. \\
\textbf{Vorbedingung}    & Alle 5 Screens implementiert. Button an GPIO\,1. \\
\textbf{Testschritte}    & 1. System starten (Tree-Screen) \newline 2. Button drücken $\rightarrow$ Overview \newline 3. Button drücken $\rightarrow$ Detail \newline 4. Button drücken $\rightarrow$ Analog \newline 5. Button drücken $\rightarrow$ Bubble \newline 6. Button drücken $\rightarrow$ zurück zu Tree (Zyklus) \newline 7. 20x schnell hintereinander drücken (Stresstest) \\
\textbf{Pass-Kriterium}  & Jeder Druck = genau 1 Wechsel. Kein Prellen. Übergang $<$300\,ms. Kein Absturz bei Schnelldrücken. Zyklus korrekt (5 Screens). \\
\textbf{Ergebnis}        & $\square$ PASS \quad $\square$ FAIL \quad Datum: \_\_\_\_\_\_ \quad Tester: \_\_\_\_\_\_ \\
\bottomrule
\end{longtable}

\begin{longtable}{p{3cm} p{10cm}}
\toprule
\rowcolor{tableheader}
\multicolumn{2}{l}{\textbf{TC-13 -- UI: Messwerte-Konsistenz alle Screens}} \\
\midrule
\endhead
\textbf{Beschreibung}   & Prüft, ob alle 5 Screens identische Messwerte anzeigen. \\
\textbf{Vorbedingung}    & Alle 5 Screens funktional. \\
\textbf{Testschritte}    & 1. Auf Overview-Screen alle Werte notieren \newline 2. Sofort zu Detail-Screen wechseln, Werte vergleichen \newline 3. Weiter zu Analog-Screen, Werte vergleichen \newline 4. Weiter zu Bubble-Screen, Werte vergleichen \newline 5. 5x wiederholen zu verschiedenen Zeiten \\
\textbf{Pass-Kriterium}  & Werte auf allen Screens identisch (Abweichung $<$1\,\% durch Timing). \\
\textbf{Ergebnis}        & $\square$ PASS \quad $\square$ FAIL \quad Datum: \_\_\_\_\_\_ \quad Tester: \_\_\_\_\_\_ \\
\bottomrule
\end{longtable}

% ============================================================
\section{Teil 5: Power Management}
% ============================================================

Display-Dimming, Aufwecken per Radar und Light Sleep.

\begin{longtable}{p{3cm} p{10cm}}
\toprule
\rowcolor{tableheader}
\multicolumn{2}{l}{\textbf{TC-14 -- Power: Display-Dimming nach Timeout}} \\
\midrule
\endhead
\textbf{Beschreibung}   & Display dimmt automatisch nach konfigurierter Zeit ohne Präsenz. \\
\textbf{Vorbedingung}    & DIM\_TIMEOUT = 45\,s. Radar aktiv. MOSFET-Steuerung funktional. \\
\textbf{Testschritte}    & 1. Präsenz auslösen (Person vor Sensor) \newline 2. Display-Helligkeit prüfen (100\,\%) \newline 3. Raum verlassen / keine Bewegung \newline 4. 50 Sekunden warten \newline 5. Helligkeit prüfen (sollte gedimmt sein \textasciitilde12\,\%) \\
\textbf{Pass-Kriterium}  & Dimming erfolgt nach 45\,s $\pm$5\,s. Helligkeit sinkt sichtbar auf Dim-Level. \\
\textbf{Ergebnis}        & $\square$ PASS \quad $\square$ FAIL \quad Datum: \_\_\_\_\_\_ \quad Tester: \_\_\_\_\_\_ \\
\bottomrule
\end{longtable}

\begin{longtable}{p{3cm} p{10cm}}
\toprule
\rowcolor{tableheader}
\multicolumn{2}{l}{\textbf{TC-15 -- Power: Display-Aufwecken per Radar}} \\
\midrule
\endhead
\textbf{Beschreibung}   & Gedimmtes Display wacht bei Präsenzerkennung sofort auf. \\
\textbf{Vorbedingung}    & Display im Dim-Modus. Radar aktiv. \\
\textbf{Testschritte}    & 1. Display dimmen lassen (keine Bewegung 45\,s) \newline 2. Person nähert sich dem Sensor ($<$2\,m) \newline 3. Reaktionszeit messen (Stoppuhr) \newline 4. Helligkeit nach Aufwecken prüfen \\
\textbf{Pass-Kriterium}  & Aufwecken innerhalb 500\,ms. Helligkeit = 100\,\%. \\
\textbf{Ergebnis}        & $\square$ PASS \quad $\square$ FAIL \quad Datum: \_\_\_\_\_\_ \quad Tester: \_\_\_\_\_\_ \\
\bottomrule
\end{longtable}

\begin{longtable}{p{3cm} p{10cm}}
\toprule
\rowcolor{tableheader}
\multicolumn{2}{l}{\textbf{TC-16 -- Power: ESP32 Light Sleep Modus}} \\
\midrule
\endhead
\textbf{Beschreibung}   & ESP32 geht nach längerer Inaktivität in Light Sleep und wacht per Radar auf. \\
\textbf{Vorbedingung}    & Light Sleep implementiert. SLEEP\_TIMEOUT = 5\,Min ohne Präsenz. \\
\textbf{Testschritte}    & 1. System normal laufen lassen \newline 2. Keine Präsenz für 5 Minuten \newline 3. Stromaufnahme messen (sollte \textasciitilde0.8\,mA im Light Sleep) \newline 4. Präsenz auslösen \newline 5. Aufwachzeit messen \newline 6. Prüfen ob alle Sensoren wieder funktionieren \\
\textbf{Pass-Kriterium}  & Light Sleep erreicht (\textasciitilde0.8\,mA). Aufwachen \textasciitilde2\,ms. Alle Sensoren nach Wake OK. \\
\textbf{Ergebnis}        & $\square$ PASS \quad $\square$ FAIL \quad Datum: \_\_\_\_\_\_ \quad Tester: \_\_\_\_\_\_ \\
\bottomrule
\end{longtable}

% ============================================================
\section{Teil 6: Systemtest}
% ============================================================

Abschließender Test aller Features im Zusammenspiel.

\begin{longtable}{p{3cm} p{10cm}}
\toprule
\rowcolor{tableheader}
\multicolumn{2}{l}{\textbf{TC-17 -- System: Komplett-Szenario Alltagsbetrieb}} \\
\midrule
\endhead
\textbf{Beschreibung}   & Simuliert einen typischen Nutzungstag mit allen Features. \\
\textbf{Vorbedingung}    & Komplett aufgebautes System. Alle Features implementiert. \\
\textbf{Testschritte}    & 1. System einschalten $\rightarrow$ Tree-Animation, volle Helligkeit \newline 2. Durch alle 5 Screens wechseln, Werte prüfen \newline 3. Raum verlassen $\rightarrow$ nach 45\,s dimmt Display \newline 4. Raum betreten $\rightarrow$ Display wacht auf \newline 5. 30 Min laufen lassen \newline 6. Werte erneut prüfen \\
\textbf{Pass-Kriterium}  & Alle Funktionen arbeiten zusammen. Keine Abstürze. Werte konsistent. \\
\textbf{Ergebnis}        & $\square$ PASS \quad $\square$ FAIL \quad Datum: \_\_\_\_\_\_ \quad Tester: \_\_\_\_\_\_ \\
\bottomrule
\end{longtable}

% ============================================================
\section{Testprotokoll}
% ============================================================

\begin{longtable}{p{1cm} p{5.5cm} p{2.5cm} p{2cm} p{2cm}}
\toprule
\rowcolor{tableheader}
\textbf{ID} & \textbf{Testname} & \textbf{Status} & \textbf{Datum} & \textbf{Tester} \\
\midrule
\endhead
TC-01 & Einzeltest: Display ST7796S           & \textcolor{pass}{PASS}    & Jan 2026 & Team \\
TC-02 & Einzeltest: AHT20 (Temp/Feuchte)      & \textcolor{pass}{PASS}    & Jan 2026 & Team \\
TC-03 & Einzeltest: SGP40 (VOC)               & \textcolor{pass}{PASS}    & Jan 2026 & Team \\
TC-04 & Einzeltest: MH-Z19C (CO\textsubscript{2}) & \textcolor{pass}{PASS} & Jan 2026 & Team \\
TC-05 & Einzeltest: PMS5003 (Feinstaub)        & \textcolor{pass}{PASS}    & Jan 2026 & Team \\
TC-06 & Einzeltest: LD2410C (Radar)            & \textcolor{pass}{PASS}    & Jan 2026 & Team \\
TC-07 & Hardware: BS170 MOSFET                 & $\square$                  &          & \\
TC-08 & Hardware: Stützkondensatoren           & \textcolor{warn}{IN ARBEIT} &        & \\
TC-09 & Integration: Alle Sensoren             & \textcolor{warn}{IN ARBEIT} &        & \\
TC-10 & UI: Screen-Darstellung (5 Screens)     & $\square$                  &          & \\
TC-11 & UI: Screen-spezifische Features        & $\square$                  &          & \\
TC-12 & UI: Button-Wechsel (5 Screens)         & $\square$                  &          & \\
TC-13 & UI: Messwerte-Konsistenz               & $\square$                  &          & \\
TC-14 & Power: Display-Dimming                 & $\square$                  &          & \\
TC-15 & Power: Aufwecken per Radar             & $\square$                  &          & \\
TC-16 & Power: ESP32 Light Sleep               & $\square$                  &          & \\
TC-17 & System: Komplett-Szenario              & $\square$                  &          & \\
\bottomrule
\end{longtable}

% ============================================================
\section{Zusammenfassung}
% ============================================================

\begin{longtable}{p{6cm} p{7cm}}
\toprule
\rowcolor{tableheader}
\textbf{Kategorie} & \textbf{Wert} \\
\midrule
\endhead
\textbf{Tests bestanden}       & \_\_\_\_ / 17 \\
\textbf{Tests nicht bestanden} & \_\_\_\_ / 17 \\
\textbf{Gesamtergebnis}        & $\square$ BESTANDEN \quad $\square$ NICHT BESTANDEN \\
\bottomrule
\end{longtable}

\vspace{1cm}
\textbf{Anmerkungen:}

\rule{\textwidth}{0.4pt}

\vspace{0.5cm}
\rule{\textwidth}{0.4pt}

\vspace{1cm}
Unterschrift: \_\_\_\_\_\_\_\_\_\_\_\_\_\_\_\_\_\_\_\_\_\_ \hfill Datum: \_\_\_\_\_\_\_\_\_\_\_\_

% ============================================================
\appendix
\section{Anhang A: Grenzwerte}
% ============================================================

\begin{longtable}{p{3cm} p{3.5cm} p{3.5cm} p{3.5cm}}
\toprule
\rowcolor{tableheader}
\textbf{Parameter} & \textbf{Gut (Grün)} & \textbf{Mäßig (Gelb)} & \textbf{Schlecht (Rot)} \\
\midrule
\endhead
\textbf{CO\textsubscript{2}} & $<$ 800\,ppm & 800--1000\,ppm & $>$ 1500\,ppm \\
\textbf{PM2.5}                & $<$ 12\,µg/m³ & 12--35\,µg/m³ & $>$ 55\,µg/m³ \\
\textbf{VOC Index}            & $<$ 100 & 100--150 & $>$ 200 \\
\bottomrule
\end{longtable}

% ============================================================
\section{Anhang B: Timeout-Konfiguration}
% ============================================================

Werte aus \texttt{power\_manager.h}:

\begin{longtable}{p{6cm} p{7cm}}
\toprule
\rowcolor{tableheader}
\textbf{Parameter} & \textbf{Wert} \\
\midrule
\endhead
\textbf{DIM\_TIMEOUT (Display dimmen)}   & 45 Sekunden ohne Präsenz \\
\textbf{SLEEP\_TIMEOUT (Light Sleep)}    & 5 Minuten ohne Präsenz \\
\textbf{BRIGHTNESS\_ACTIVE}              & 255 (100\,\%) \\
\textbf{BRIGHTNESS\_DIMMED}              & 30 (\textasciitilde12\,\%) \\
\textbf{BRIGHTNESS\_OFF}                 & 0 (Display aus) \\
\bottomrule
\end{longtable}

% ============================================================
\section{Anhang C: Kritische Hardware-Hinweise}
% ============================================================

\textbf{Stützkondensatoren (220\,µF 25V):} ZWINGEND erforderlich an VCC/GND von MH-Z19C und LD2410C! Ohne diese stürzt das System bei Stromspitzen ab.

\textbf{Level Shifter (LD2410C):} TX-Ausgang des Radar-Sensors MUSS über Level Shifter an GPIO\,7 laufen. Direkte 5V-Signale zerstören den ESP32!

\textbf{BS170 MOSFET:} Gate-Widerstand 1\,k$\Omega$ zwischen GPIO\,3 und Gate. Logic-Level MOSFET, schaltet bei 3.3V durch.

\end{document}
